%%=============================================================================
%% Samenvatting
%%=============================================================================

% TODO: De "abstract" of samenvatting is een kernachtige (~ 1 blz. voor een thesis) synthese van het document.
%
% Een goede abstract biedt een kernachtig antwoord op volgende vragen:
%
% 1. Waarover gaat de bachelorproef?
% 2. Waarom heb je er over geschreven?
% 3. Hoe heb je het onderzoek uitgevoerd?
% 4. Wat waren de resultaten? Wat blijkt uit je onderzoek?
% 5. Wat betekenen je resultaten? Wat is de relevantie voor het werkveld?
%
% Daarom bestaat een abstract uit volgende componenten:
%
% - inleiding + kaderen thema
% - probleemstelling
% - (centrale) onderzoeksvraag
% - onderzoeksdoelstelling
% - methodologie
% - resultaten (beperk tot de belangrijkste, relevant voor de onderzoeksvraag)
% - conclusies, aanbevelingen, beperkingen
%
% LET OP! Een samenvatting is GEEN voorwoord!

%%---------- Samenvatting -----------------------------------------------------
% De samenvatting in de hoofdtaal van het document

\chapter*{\IfLanguageName{dutch}{Samenvatting}{Abstract}}

De Europese Corporate Sustainability Reporting Directive (CSRD) verplicht productiebedrijven tot gedetailleerde rapportage over milieu-, sociale en governanceaspecten conform de European Sustainability Reporting Standards (ESRS). Voor organisaties zoals Turtle Srl brengt dit een aanzienlijke administratieve last met zich mee: Environmental, Social en Governance (ESG)-informatie is verspreid over heterogene documenten in meerdere formaten en talen, terwijl externe partijen via vragenlijsten zoals EcoVadis en de Voluntary Sustainability Reporting Standard (VSME) een transparante verantwoording van duurzaamheidsprestaties verwachten. Standaard Large Language Models (LLM's) bieden onvoldoende soelaas: ze kennen geen interne bedrijfsdocumenten, zijn vatbaar voor hallucinaties en missen de traceerbare bronvermelding die CSRD-auditing vereist.

Deze bachelorproef onderzoekt hoe Retrieval-Augmented Generation (RAG) kan worden ingezet als oplossing voor deze tekortkomingen. De centrale onderzoeksvraag luidt: \textit{Welke combinatie van chunkingstrategie, embeddingmodel en Large Language Model (LLM) resulteert in een optimaal RAG-systeem voor Turtle Srl's Proof of Concept (PoC)?} Het beoogde resultaat is een vergelijkende studie en een technisch adviesrapport dat Turtle Srl begeleidt bij de keuze van de optimale RAG-pipelineconfiguratie, met als kritische succesfactoren feitelijke betrouwbaarheid (faithfulness), juridische traceerbaarheid (source attribution) en contextvaliditeit (context recall).

Het onderzoek volgt een sequentieel framework van zeven fasen. Na een uitgebreid literatuuronderzoek en een requirementsanalyse via de MoSCoW-methode, werden drie chunkingstrategieën (\texttt{RecursiveCharacterSplitter}, \texttt{FixedSizeTokenSplitter} en \texttt{SemanticEmbeddingChunker}), drie embeddingmodellen (\texttt{multilingual-e5-large-instruct}, \texttt{BAAI/bge-m3} en \texttt{Snowflake Arctic Embed}) en drie LLM's (\texttt{Gemma-3-27b-it}, \texttt{Mistral-Small-2603} en \texttt{Llama-3.3-70B-Instruct}) geselecteerd via een longlist-shortlistprocedure. De experimentele dataset bestond uit interne ESG-documenten van Turtle Srl en een handmatig geverifieerde golden dataset van 20 vraag-antwoordparen. Met Haystack~2.x en het HyPSTER-configuratieframework werden alle 27 unieke componentcombinaties ($3 \times 3 \times 3$) systematisch geëvalueerd over 540 runs. De performantie werd gemeten via vier Retrieval Augmented Generation Assessment (RAGAS)-metrics — faithfulness, answer relevancy, context precision en context recall — aangevuld met een domeinspecifieke source attribution-score die de correctheid van bestandsnaam en paginanummer beoordeelt. Statistische vergelijkingen werden uitgevoerd via de Kruskal-Wallis-toets met Dunn/Holm post-hoc correctie, en gevalideerd via een menselijke validatiestudie.

De resultaten onthullen een duidelijke hiërarchie in componenteffecten. Het embeddingmodel is de dominante kwaliteitshefboom: \texttt{multilingual-e5-large-instruct} behaalt de hoogste context recall ($\mu = 0{,}792$) en vermijdt de source attribution-penalisering die \texttt{bge-m3} treft door zijn hybride Reciprocal Rank Fusion (RRF)-retrieval. De chunkingstrategie is doorslaggevend voor source attribution: de \texttt{RecursiveCharacterSplitter} is als enige strategie in staat SA\,$\geq$\,40\,\% te halen (46{,}7\,\%, $p_{\text{holm}} < 0{,}001$), doordat zijn scheidingstekenhiërarchie paginagrensmarkeerders structureel intact laat. Opmerkelijk is dat het LLM geen statistisch significant effect vertoont op faithfulness noch op answer relevancy ($H = 5{,}78$, $p = 0{,}056$ resp.\ $H = 5{,}00$, $p = 0{,}082$). De combinatie \textbf{\texttt{RecursiveCharacterSplitter} + \texttt{multilingual-e5-large-instruct} + \texttt{Gemma-3-27b-it}} wordt formeel aanbevolen als optimale configuratie voor de PoC van Turtle Srl.

Het meest fundamentele inzicht voor de implementatiepraktijk is dat optimalisatie-inspanningen bij RAG-pipelines voor ESG-compliance beter gericht zijn op de upstream vectorisatie- en chunkfase dan op de selectie van een krachtigere generatieve engine. Op basis van de experimentele resultaten werd een metriekgestuurde beslissingsboom opgesteld die toekomstige implementatoren begeleidt bij de configuratiekeuze afhankelijk van hun prioriteit: juridische traceerbaarheid, antwoordrelevantie of contextdekking. De generaliseerbaarheid van de resultaten is beperkt tot vergelijkbare meertalige ESG-documentstromen; validatie op andere domeinen of met uitgebreidere golden datasets vormt een logische vervolgstap.
